% Chapter 2 outline. Add prose only after collecting and evaluating reliable
% scholarly, government, standards, and official technical sources. Do not
% cite the selected-capstone slides as scholarly evidence.

\newpage
\section{Review of Related Literature}
\label{Review of Related Literature}
% Introduce the purpose, scope, organization, inclusion criteria, publication
% period, and source types used in the review.
% End the chapter with a synthesis that connects the literature to the actual
% research gap and proposed contribution.

\subsection{Flood Evacuation Planning and Decision Support}
% Synthesize concepts relevant to the problem rather than listing definitions.
% Establish how flood conditions affect accessibility, evacuation demand,
% shelter selection, travel time, and allocation of limited resources.

\subsubsection{Flood Risk, Road Accessibility, and Evacuation Centers}
% Review evidence on:
% - flood hazards and their effects on transportation networks;
% - road and bridge disruption during evacuation;
% - evacuation-center location, accessibility, capacity, and suitability;
% - population and resource information used by local disaster planners;
% - drainage and terrain data only after their role in this study is decided.

\subsubsection{Scenario-Based Simulation, Routing, and Optimization}
% Review evidence on:
% - what-if and scenario-based disaster planning;
% - decision-support systems for emergency management;
% - graph-based evacuation routing under changing road conditions;
% - capacity-constrained shelter assignment and resource allocation;
% - evacuation-time estimation;
% - hydrologic/hydraulic flood modeling only if it enters the final scope.
% Define “simulation” and “optimal” consistently with the selected method.

\subsection{Related Systems}
% Compare existing local and international flood/evacuation tools.
% For each system record:
% - intended users and geographic scope;
% - inputs and data sources;
% - simulation and decision methods;
% - route, shelter, time, and resource outputs;
% - scenario-comparison capability;
% - deployment model and limitations.
% Use a comparison table after verified systems are selected. The synthesis
% should show the specific gap addressed by this project without unsupported
% novelty claims.

\subsection{Technical Background}
% Explain only technologies that are part of the approved design:
% - GIS concepts, coordinate reference systems, and spatial data quality;
% - road networks represented as graphs;
% - drainage, waterways, terrain, hazards, population, and center layers;
% - GeoJSON or other selected formats and the validation pipeline;
% - routing, allocation, travel-time, and scenario algorithms;
% - system architecture, database, API, interface, and deployment environment;
% - privacy, access control, backup, and data-update responsibilities.
% OpenStreetMap, QGIS, Python, and FastAPI remain tentative until approved.

\subsection{Conceptual Framework}
% Present and explain the final framework using a figure based on:
% Input -> Processing -> Output -> User Decision -> Measurable Benefit.
% Inputs should cover only confirmed layers and scenario variables.
% Processing must name the approved algorithms and constraints.
% Outputs must map to the objectives and evaluation measures.
% Source foundation: knowledge/00-capstone-definition.md, Section 7 and the
% Definition Integrity Test.
